\cvsection{Experience}
\newcommand{\quotes}[1]{``#1''}

\begin{cventries}

  %---------------------------------------------------------
  \cventry
  {Frontend Developer}
  {Publicis Global Delivery (PGD) — assigned to Digitas}
  {Colombia (Remote)}
  {Jan. 2025 -- Present}
  {
    \begin{cvitems}
      \item {Collaborated with Product \& Design to translate requirements into \textbf{accessible, maintainable UI components} using \textbf{Next.js + TypeScript}.}
      \item {Shipped production features and stability fixes to an internal \textbf{Nuxt.js} web app, so I’m comfortable working in a Vue-based stack with TypeScript.}
      \item {Optimized builds and integrated CI/CD using \textbf{GitHub Actions}, similar to PayPay Card’s static-deploy model.}
      \item {Improved runtime performance and error visibility through monitoring tools (Dynatrace / Sentry-like) to keep UIs stable.}
    \end{cvitems}
  }

  %---------------------------------------------------------
  \cventry
  {React Native Developer}
  {Publicis Global Delivery (PGD) — assigned to Razorfish / Wegmans}
  {Colombia (Remote \textemdash{} New York)}
  {Dec. 2023 -- Dec. 2024}
  {
    \begin{cvitems}
      \item {Delivered \textbf{production React Native features} at scale using \textbf{TypeScript}, \textbf{Zustand}, and \textbf{Algolia}, increasing in-app engagement by \textbf{30\%}.}
      \item {Integrated monitoring (Dynatrace) and \textbf{Jest} unit tests, reducing critical bugs by \textbf{25\%}.}
      \item {Automated \textbf{CI/CD} with GitHub Actions for mobile builds, supporting iterative delivery with small, safe releases.}
      \item {Implemented responsive, accessible UI for touch and keyboard input, focusing on performance and predictable state.}
    \end{cvitems}
  }

  %---------------------------------------------------------
  \cventry
  {Full Stack Developer (Self-employed)}
  {Grisú}
  {Colombia (Remote)}
  {Jul.\ 2023 -- Present}
  {
    \begin{cvitems}
      \item {Built full-stack web apps with \textbf{React + NestJS} and deployed to \textbf{Azure/AWS} with secure CI/CD pipelines.}
      \item {Worked directly with stakeholders (including public-sector programs in Curaçao) to gather requirements and iterate fast.}
    \end{cvitems}
  }

  %---------------------------------------------------------
  \cventry
  {Full Stack Developer (Part-time)}
  {ArcDesign}
  {Spain (Remote)}
  {Jul.\ 2023 -- Mar.\ 2025}
  {
    \begin{cvitems}
      \item {Implemented a geolocation algorithm using \textbf{hashmaps + geohashes}, improving scalability and accuracy.}
      \item {Integrated that logic in \textbf{React Native} for fast, near real-time UX.}
      \item {Deployed React-based sites to AWS/Hostinger.}
    \end{cvitems}
  }

  %---------------------------------------------------------
  \cventry
  {Problem Solver / Full Stack Developer}
  {Design Systems Inno}
  {Medellín, Colombia (Hybrid)}
  {Dec.\ 2022 -- Jun.\ 2023}
  {
    \begin{cvitems}
      \item {Designed internal APIs using \textbf{Go Fiber}, \textbf{Python FastAPI} and \textbf{MongoDB} with \textbf{RabbitMQ} for async tasks.}
      \item {Created React dashboards (Redux + Saga) and \textbf{optimized React and Angular builds}, reducing bundle sizes from 3.46 MB to \textbf{~886 KB}.}
      \item {Designed internal APIs with \textbf{Go Fiber}, \textbf{FastAPI} and \textbf{MongoDB} using \textbf{RabbitMQ} for async processing.}
      \item \item {Contributed small \textbf{Vue.js} views/components in mixed-frontends projects, so I’m familiar with the Vue SFC workflow.}
    \end{cvitems}
  }

  %---------------------------------------------------------
  \cventry
  {Cyber Security Analyst / Red Team}
  {Geta Club Play \& Aligo}
  {Colombia}
  {Sep.\ 2021 -- Oct.\ 2022}
  {
    \begin{cvitems}
      \item {Implemented DevSecOps pipelines using \textbf{OWASP ZAP} and \textbf{SonarQube}, reducing vulnerabilities by \textbf{50\%} in 4 months.}
      \item {Reverse engineered APKs to detect sensitive info exposure and worked with devs to mitigate.}
      \item {Wrote automation scripts in Python and contributed to secure coding practices.}
    \end{cvitems}
  }

\end{cventries}
